\section{Theory of change}
Conditional on the underlying model of ontology being true, the
techniques proposed here would be a substantial interpretability
advance. The agenda also has a case for counterfactual impact: its
theoretical basis consists of nonobvious applications of obscure
techniques from niche fields of bioinformatics, and similar research
is unlikely to be done soon otherwise.

\subsection{What a solution would look like}
Mechinterp plausibly offers the best path to a robust alignment
solution despite
\href{https://www.lesswrong.com/posts/tEPHGZAb63dfq2v8n/how-useful-is-mechanistic-interpretability}{criticisms}.
The case rests on
\begin{enumerate}
    \item \textbf{decomposability} Mechinterp allows alignment to be
    broken into manageable subproblems.
    \item \textbf{parallelizability} Decomposable subproblems can be
    worked on independently.
    \item \textbf{robustness} Mechinterp offers direct causal
    explanations of model behavior.
\end{enumerate}
The most straightforward mechinterp alignment solutions are
\begin{enumerate}
    \item \href{https://www.alignment.org/blog/arcs-first-technical-report-eliciting-latent-knowledge/}{eliciting latent knowledge (ELK)}
    Directly inferring a model's ontology.
    \label{it:elk}
    \item \href{https://www.alignment.org/blog/mechanistic-anomaly-detection-and-elk/}{mechanistic anomaly detection (MAD)}
    Detecting when a model has ``unusual'' thoughts.
    \label{it:mad}
\end{enumerate}
Both of these can be considered
\href{https://arbital.greaterwrong.com/p/ontology_identification/}{ontology identification}
subproblems. MAD requires mapping between the model's activations and
the model's ontology. ELK requires mapping between the model's ontology
and human ontology. These are hard but not intractable problems.
Recently, \href{https://devinterp.com/}{developmental interpretability}
(devinterp) techniques have enabled identification of ``developmental
stages'' corresponding to learning specific
tasks\cite{hoogland2024developmental}. This is strong evidence that
learning consists primarily of the network
grokking\cite{nanda2023progress} discrete concepts. While these results
are promising for MAD, locating these emergent capabilities in the
model's activations remains an open problem.

Interpretability glosses broadly as ``let humans understand AI''. The
alternative is to get AI to understand humans better, and that scales
badly: making the model smarter leads to subtler and subtler issues
with larger and larger effects. And what we're doing is not just
letting humans understand AI. Because an LLM reflects humanity from its
training data, high quality conceptual interpretability provides
machine-comprehensible mathematical representations of every concept
that humanity cares about. We work on interp for the usual
corrigibility reasons, and also because feature representation is one
of the most challenging parts of outer alignment. This attacks both.

\subsection{Macro-scale strategy}
Broadly, the goals are to
\begin{enumerate}
  \item formalize the notion of a model's latent concept space
  \item distill a model's latent concept space
  \item enforce decomposability of latent concept space
  \item enforce robustness of existing concepts as the dimensionality of
  the latent concept space increases
\end{enumerate}
A prerequisite for all of these is the ability to express a
\href{https://plato.stanford.edu/entries/ontological-commitment/}{metaontology}
which can identify latent concepts in activation space.


